\documentclass[aspectratio=169, 11pt]{beamer}
% ============================================================================
% Preambulo compartido para presentaciones Beamer
% Estadistica 2026-II - Pontificia Universidad Javeriana
% Incluir con: % ============================================================================
% Preambulo compartido para presentaciones Beamer
% Estadistica 2026-II - Pontificia Universidad Javeriana
% Incluir con: % ============================================================================
% Preambulo compartido para presentaciones Beamer
% Estadistica 2026-II - Pontificia Universidad Javeriana
% Incluir con: \input{../preambulo_beamer.tex} despues de \documentclass
% ============================================================================

\usepackage[T1]{fontenc}
\usepackage[utf8]{inputenc}
\usepackage[spanish]{babel}
\usepackage{amsmath, amssymb, amsfonts}
\usepackage{booktabs}
\usepackage{graphicx}
\usepackage{tikz}
\usetikzlibrary{shapes.geometric, positioning, arrows.meta, calc}
\usepackage{pgfplots}
\pgfplotsset{compat=1.18}
\usepackage{listings}
\usepackage{xcolor}
\usepackage{hyperref}
\usepackage{multicol}

% --- Tema Beamer (minimalista) ---
\usetheme{default}
\usecolortheme{default}
\usefonttheme{default}
\setbeamertemplate{navigation symbols}{}
\setbeamertemplate{caption}[numbered]
\setbeamertemplate{footline}{%
  \hspace{1em}{\tiny\url{https://github.com/polanco-jaime/Curso-Estadistica-2026-3}}\hfill\usebeamerfont{footline}\insertframenumber/\inserttotalframenumber\hspace{1em}\vspace{0.5em}%
}
\setbeamertemplate{itemize items}[circle]
\setbeamertemplate{enumerate items}[default]

% --- Colores institucionales (sutiles) ---
\definecolor{javeriana}{RGB}{0, 62, 105}
\definecolor{javerianalight}{RGB}{0, 105, 170}
\definecolor{codigobg}{RGB}{245, 245, 245}
\definecolor{codigofg}{RGB}{40, 40, 40}
\definecolor{comentario}{RGB}{100, 100, 100}
\definecolor{keyword}{RGB}{0, 100, 150}
\definecolor{string}{RGB}{180, 60, 30}

\setbeamercolor{title}{fg=javeriana}
\setbeamercolor{frametitle}{fg=javeriana}
\setbeamercolor{structure}{fg=javeriana}
\setbeamercolor{block title}{fg=white, bg=javeriana!75}
\setbeamercolor{block body}{bg=javeriana!5}
\setbeamercolor{block title alerted}{fg=white, bg=red!70!black}
\setbeamercolor{block body alerted}{bg=red!5}
\setbeamercolor{block title example}{fg=white, bg=green!40!black}
\setbeamercolor{block body example}{bg=green!5}
\setbeamercolor{item}{fg=javeriana}

\setbeamerfont{title}{series=\bfseries, size=\Large}
\setbeamerfont{frametitle}{series=\bfseries}

% --- Configuracion de listings para R ---
\lstset{
  language=R,
  basicstyle=\ttfamily\footnotesize,
  backgroundcolor=\color{codigobg},
  commentstyle=\color{comentario}\itshape,
  keywordstyle=\color{keyword}\bfseries,
  stringstyle=\color{string},
  numbers=none,
  frame=single,
  rulecolor=\color{javeriana!30},
  breaklines=true,
  showstringspaces=false,
  columns=flexible,
  morekeywords={library, ggplot, aes, geom_histogram, geom_boxplot,
    geom_point, geom_smooth, geom_density, geom_bar, geom_line,
    geom_vline, geom_hline, geom_errorbar, labs, theme_minimal,
    facet_wrap, scale_fill_manual, scale_color_manual,
    summarise, group_by, filter, mutate, select, arrange,
    read.csv, head, str, summary, mean, median, sd, var, cor,
    lm, t.test, chisq.test, wilcox.test, kruskal.test, aov,
    pnorm, qnorm, dnorm, rnorm, qt, pt, qchisq, pchisq, qf, pf,
    confint, coeftest, vcovHC, bptest, vif, cooks.distance,
    TukeyHSD, table, prop.table}
}

% --- Comandos personalizados ---
\newcommand{\R}{\texttt{R}}
\newcommand{\code}[1]{\texttt{\footnotesize #1}}
\newcommand{\variable}[1]{\texttt{#1}}
\newcommand{\paquete}[1]{\texttt{#1}}

% --- Informacion del curso ---
\institute{Pontificia Universidad Javeriana\\
  Facultad de Ciencias Econ\'omicas y Administrativas}
\date{2026-II}
 despues de \documentclass
% ============================================================================

\usepackage[T1]{fontenc}
\usepackage[utf8]{inputenc}
\usepackage[spanish]{babel}
\usepackage{amsmath, amssymb, amsfonts}
\usepackage{booktabs}
\usepackage{graphicx}
\usepackage{tikz}
\usetikzlibrary{shapes.geometric, positioning, arrows.meta, calc}
\usepackage{pgfplots}
\pgfplotsset{compat=1.18}
\usepackage{listings}
\usepackage{xcolor}
\usepackage{hyperref}
\usepackage{multicol}

% --- Tema Beamer (minimalista) ---
\usetheme{default}
\usecolortheme{default}
\usefonttheme{default}
\setbeamertemplate{navigation symbols}{}
\setbeamertemplate{caption}[numbered]
\setbeamertemplate{footline}{%
  \hspace{1em}{\tiny\url{https://github.com/polanco-jaime/Curso-Estadistica-2026-3}}\hfill\usebeamerfont{footline}\insertframenumber/\inserttotalframenumber\hspace{1em}\vspace{0.5em}%
}
\setbeamertemplate{itemize items}[circle]
\setbeamertemplate{enumerate items}[default]

% --- Colores institucionales (sutiles) ---
\definecolor{javeriana}{RGB}{0, 62, 105}
\definecolor{javerianalight}{RGB}{0, 105, 170}
\definecolor{codigobg}{RGB}{245, 245, 245}
\definecolor{codigofg}{RGB}{40, 40, 40}
\definecolor{comentario}{RGB}{100, 100, 100}
\definecolor{keyword}{RGB}{0, 100, 150}
\definecolor{string}{RGB}{180, 60, 30}

\setbeamercolor{title}{fg=javeriana}
\setbeamercolor{frametitle}{fg=javeriana}
\setbeamercolor{structure}{fg=javeriana}
\setbeamercolor{block title}{fg=white, bg=javeriana!75}
\setbeamercolor{block body}{bg=javeriana!5}
\setbeamercolor{block title alerted}{fg=white, bg=red!70!black}
\setbeamercolor{block body alerted}{bg=red!5}
\setbeamercolor{block title example}{fg=white, bg=green!40!black}
\setbeamercolor{block body example}{bg=green!5}
\setbeamercolor{item}{fg=javeriana}

\setbeamerfont{title}{series=\bfseries, size=\Large}
\setbeamerfont{frametitle}{series=\bfseries}

% --- Configuracion de listings para R ---
\lstset{
  language=R,
  basicstyle=\ttfamily\footnotesize,
  backgroundcolor=\color{codigobg},
  commentstyle=\color{comentario}\itshape,
  keywordstyle=\color{keyword}\bfseries,
  stringstyle=\color{string},
  numbers=none,
  frame=single,
  rulecolor=\color{javeriana!30},
  breaklines=true,
  showstringspaces=false,
  columns=flexible,
  morekeywords={library, ggplot, aes, geom_histogram, geom_boxplot,
    geom_point, geom_smooth, geom_density, geom_bar, geom_line,
    geom_vline, geom_hline, geom_errorbar, labs, theme_minimal,
    facet_wrap, scale_fill_manual, scale_color_manual,
    summarise, group_by, filter, mutate, select, arrange,
    read.csv, head, str, summary, mean, median, sd, var, cor,
    lm, t.test, chisq.test, wilcox.test, kruskal.test, aov,
    pnorm, qnorm, dnorm, rnorm, qt, pt, qchisq, pchisq, qf, pf,
    confint, coeftest, vcovHC, bptest, vif, cooks.distance,
    TukeyHSD, table, prop.table}
}

% --- Comandos personalizados ---
\newcommand{\R}{\texttt{R}}
\newcommand{\code}[1]{\texttt{\footnotesize #1}}
\newcommand{\variable}[1]{\texttt{#1}}
\newcommand{\paquete}[1]{\texttt{#1}}

% --- Informacion del curso ---
\institute{Pontificia Universidad Javeriana\\
  Facultad de Ciencias Econ\'omicas y Administrativas}
\date{2026-II}
 despues de \documentclass
% ============================================================================

\usepackage[T1]{fontenc}
\usepackage[utf8]{inputenc}
\usepackage[spanish]{babel}
\usepackage{amsmath, amssymb, amsfonts}
\usepackage{booktabs}
\usepackage{graphicx}
\usepackage{tikz}
\usetikzlibrary{shapes.geometric, positioning, arrows.meta, calc}
\usepackage{pgfplots}
\pgfplotsset{compat=1.18}
\usepackage{listings}
\usepackage{xcolor}
\usepackage{hyperref}
\usepackage{multicol}

% --- Tema Beamer (minimalista) ---
\usetheme{default}
\usecolortheme{default}
\usefonttheme{default}
\setbeamertemplate{navigation symbols}{}
\setbeamertemplate{caption}[numbered]
\setbeamertemplate{footline}{%
  \hspace{1em}{\tiny\url{https://github.com/polanco-jaime/Curso-Estadistica-2026-3}}\hfill\usebeamerfont{footline}\insertframenumber/\inserttotalframenumber\hspace{1em}\vspace{0.5em}%
}
\setbeamertemplate{itemize items}[circle]
\setbeamertemplate{enumerate items}[default]

% --- Colores institucionales (sutiles) ---
\definecolor{javeriana}{RGB}{0, 62, 105}
\definecolor{javerianalight}{RGB}{0, 105, 170}
\definecolor{codigobg}{RGB}{245, 245, 245}
\definecolor{codigofg}{RGB}{40, 40, 40}
\definecolor{comentario}{RGB}{100, 100, 100}
\definecolor{keyword}{RGB}{0, 100, 150}
\definecolor{string}{RGB}{180, 60, 30}

\setbeamercolor{title}{fg=javeriana}
\setbeamercolor{frametitle}{fg=javeriana}
\setbeamercolor{structure}{fg=javeriana}
\setbeamercolor{block title}{fg=white, bg=javeriana!75}
\setbeamercolor{block body}{bg=javeriana!5}
\setbeamercolor{block title alerted}{fg=white, bg=red!70!black}
\setbeamercolor{block body alerted}{bg=red!5}
\setbeamercolor{block title example}{fg=white, bg=green!40!black}
\setbeamercolor{block body example}{bg=green!5}
\setbeamercolor{item}{fg=javeriana}

\setbeamerfont{title}{series=\bfseries, size=\Large}
\setbeamerfont{frametitle}{series=\bfseries}

% --- Configuracion de listings para R ---
\lstset{
  language=R,
  basicstyle=\ttfamily\footnotesize,
  backgroundcolor=\color{codigobg},
  commentstyle=\color{comentario}\itshape,
  keywordstyle=\color{keyword}\bfseries,
  stringstyle=\color{string},
  numbers=none,
  frame=single,
  rulecolor=\color{javeriana!30},
  breaklines=true,
  showstringspaces=false,
  columns=flexible,
  morekeywords={library, ggplot, aes, geom_histogram, geom_boxplot,
    geom_point, geom_smooth, geom_density, geom_bar, geom_line,
    geom_vline, geom_hline, geom_errorbar, labs, theme_minimal,
    facet_wrap, scale_fill_manual, scale_color_manual,
    summarise, group_by, filter, mutate, select, arrange,
    read.csv, head, str, summary, mean, median, sd, var, cor,
    lm, t.test, chisq.test, wilcox.test, kruskal.test, aov,
    pnorm, qnorm, dnorm, rnorm, qt, pt, qchisq, pchisq, qf, pf,
    confint, coeftest, vcovHC, bptest, vif, cooks.distance,
    TukeyHSD, table, prop.table}
}

% --- Comandos personalizados ---
\newcommand{\R}{\texttt{R}}
\newcommand{\code}[1]{\texttt{\footnotesize #1}}
\newcommand{\variable}[1]{\texttt{#1}}
\newcommand{\paquete}[1]{\texttt{#1}}

% --- Informacion del curso ---
\institute{Pontificia Universidad Javeriana\\
  Facultad de Ciencias Econ\'omicas y Administrativas}
\date{2026-II}

\title{Sesión 9: Regresión Lineal Simple (OLS)}
\subtitle{El método de mínimos cuadrados ordinarios}
\author{Prof.\ Jaime Polanco Jim\'enez}
\date{Vie 23 oct 2026}

\begin{document}

\begin{frame}
\titlepage
\end{frame}

\begin{frame}{Agenda}
\tableofcontents
\end{frame}

% ============================================================================
\section{Introducción a la regresión}
% ============================================================================

\begin{frame}{Motivación: más allá de la correlación}
\begin{block}{Hasta ahora}
Hemos estudiado la \textbf{correlación} $r$ entre dos variables numéricas, que mide la fuerza y dirección de la relación lineal.
\end{block}

\vspace{0.2cm}
\textbf{Limitaciones de la correlación:}
\begin{itemize}
  \item Solo mide asociación, no permite predecir
  \item Es simétrica: $\text{cor}(X,Y) = \text{cor}(Y,X)$
  \item No distingue variable dependiente de independiente
\end{itemize}

\vspace{0.3cm}
\begin{exampleblock}{Pregunta de investigación}
¿Podemos \textbf{predecir} el desempeño en inglés en Saber Pro (MOD\_INGLES\_PUNT) a partir del puntaje de inglés en Saber 11 (INGLES\_SABER11\_PUNT)?

Necesitamos un \textbf{modelo de regresión}.
\end{exampleblock}
\end{frame}

\begin{frame}{El modelo de regresión lineal simple}
\begin{block}{Definición}
El modelo de \textbf{regresión lineal simple} relaciona una variable dependiente $Y$ con una variable independiente $X$ mediante una línea recta:
\[
Y_i = \beta_0 + \beta_1 X_i + \varepsilon_i
\]
donde:
\begin{itemize}
  \item $Y_i$ = valor observado de la variable dependiente (respuesta)
  \item $X_i$ = valor de la variable independiente (predictor, regressor)
  \item $\beta_0$ = intercepto (valor de $Y$ cuando $X = 0$)
  \item $\beta_1$ = pendiente (cambio en $Y$ por cada unidad de cambio en $X$)
  \item $\varepsilon_i$ = error aleatorio (desviación del modelo)
\end{itemize}
\end{block}

\vspace{0.2cm}
\textbf{Interpretación:}
\begin{itemize}
  \item $\beta_0$ y $\beta_1$ son \textbf{parámetros poblacionales} (desconocidos)
  \item Los estimaremos con $\hat{\beta}_0$ y $\hat{\beta}_1$ usando los datos
\end{itemize}
\end{frame}

\begin{frame}{Ecuación estimada y valores ajustados}
La \textbf{ecuación estimada} (fitted line) es:
\[
\hat{Y}_i = \hat{\beta}_0 + \hat{\beta}_1 X_i
\]

\vspace{0.2cm}
\begin{itemize}
  \item $\hat{Y}_i$ = valor \textbf{predicho} (ajustado) de $Y$ para la observación $i$
  \item $e_i = Y_i - \hat{Y}_i$ = \textbf{residuo} (error de predicción)
\end{itemize}

\vspace{0.3cm}
\begin{exampleblock}{Ejemplo}
Si $\hat{\beta}_0 = 45$ y $\hat{\beta}_1 = 1.8$, la ecuación es:
\[
\widehat{\text{MOD\_INGLES\_PUNT}} = 45 + 1.8 \times \text{INGLES\_SABER11\_PUNT}
\]

Para un estudiante con puntaje 60 en Saber 11:
\[
\hat{Y} = 45 + 1.8 \times 60 = 45 + 108 = 153
\]

Si su puntaje real en Saber Pro es 158, el residuo es $e = 158 - 153 = 5$.
\end{exampleblock}
\end{frame}

\begin{frame}[fragile]{Visualización: scatter plot + línea de regresión}
\begin{lstlisting}
# Cargar datos ICFES para Negocios Internacionales
icfes_ni <- read.csv("datos/icfes_ni_2025.csv")

# Scatter plot con linea de regresion
library(ggplot2)

ggplot(icfes_ni, aes(x = INGLES_SABER11_PUNT,
                     y = MOD_INGLES_PUNT)) +
  geom_point(alpha = 0.3, color = "steelblue") +
  geom_smooth(method = "lm", se = TRUE, color = "red",
              fill = "pink", alpha = 0.2) +
  labs(title = "Relacion entre ingles Saber 11 y Saber Pro",
       x = "Puntaje ingles Saber 11",
       y = "Puntaje ingles Saber Pro (MOD_INGLES_PUNT)") +
  theme_minimal()
\end{lstlisting}

\vspace{0.2cm}
La banda sombreada representa el intervalo de confianza del 95\% para la línea de regresión.
\end{frame}

\begin{frame}{¿Qué línea elegir?}
Hay infinitas líneas que pasan por la nube de puntos. ¿Cuál es la \textbf{mejor}?

\vspace{0.2cm}
\begin{block}{Criterio: minimizar los errores}
Queremos que los residuos $e_i = Y_i - \hat{Y}_i$ sean lo más pequeños posible en conjunto.
\end{block}

\vspace{0.2cm}
\textbf{Opciones:}
\begin{itemize}
  \item Minimizar $\sum |e_i|$ (suma de errores absolutos)
  \item Minimizar $\max |e_i|$ (error máximo)
  \item Minimizar $\sum e_i^2$ (suma de errores al cuadrado) $\leftarrow$ \textbf{OLS}
\end{itemize}

\vspace{0.3cm}
\begin{alertblock}{Método de Mínimos Cuadrados Ordinarios (OLS)}
El método \textbf{OLS} (Ordinary Least Squares) elige $\hat{\beta}_0$ y $\hat{\beta}_1$ que minimizan:
\[
\text{SSE} = \sum_{i=1}^{n} e_i^2 = \sum_{i=1}^{n} (Y_i - \hat{Y}_i)^2 = \sum_{i=1}^{n} (Y_i - \hat{\beta}_0 - \hat{\beta}_1 X_i)^2
\]
\end{alertblock}
\end{frame}

% ============================================================================
\section{Método de mínimos cuadrados (OLS)}
% ============================================================================

\begin{frame}{Derivación: ecuaciones normales}
Para minimizar $\text{SSE} = \sum (Y_i - \beta_0 - \beta_1 X_i)^2$, tomamos las derivadas parciales respecto a $\beta_0$ y $\beta_1$ e igualamos a cero:

\vspace{0.2cm}
\textbf{Derivada respecto a $\beta_0$:}
\[
\frac{\partial \text{SSE}}{\partial \beta_0} = -2 \sum (Y_i - \beta_0 - \beta_1 X_i) = 0
\]
\[
\Rightarrow \sum Y_i = n \beta_0 + \beta_1 \sum X_i
\]

\vspace{0.2cm}
\textbf{Derivada respecto a $\beta_1$:}
\[
\frac{\partial \text{SSE}}{\partial \beta_1} = -2 \sum X_i(Y_i - \beta_0 - \beta_1 X_i) = 0
\]
\[
\Rightarrow \sum X_i Y_i = \beta_0 \sum X_i + \beta_1 \sum X_i^2
\]

\vspace{0.2cm}
Estas son las \textbf{ecuaciones normales}. Resolviéndolas obtenemos los estimadores OLS.
\end{frame}

\begin{frame}{Fórmulas de los estimadores OLS}
\begin{block}{Estimador de la pendiente}
\[
\hat{\beta}_1 = \frac{\sum_{i=1}^{n}(X_i - \bar{X})(Y_i - \bar{Y})}{\sum_{i=1}^{n}(X_i - \bar{X})^2} = \frac{S_{XY}}{S_{XX}}
\]
donde:
\begin{itemize}
  \item $S_{XY} = \sum(X_i - \bar{X})(Y_i - \bar{Y})$ = suma de productos cruzados
  \item $S_{XX} = \sum(X_i - \bar{X})^2$ = suma de cuadrados de $X$
\end{itemize}
\end{block}

\begin{block}{Estimador del intercepto}
\[
\hat{\beta}_0 = \bar{Y} - \hat{\beta}_1 \bar{X}
\]
\end{block}

\vspace{0.2cm}
\textbf{Propiedad importante:} La línea OLS \textbf{siempre} pasa por el punto $(\bar{X}, \bar{Y})$.
\end{frame}

\begin{frame}{Relación con la correlación}
El estimador OLS de la pendiente se puede escribir como:
\[
\hat{\beta}_1 = r \times \frac{s_Y}{s_X}
\]
donde:
\begin{itemize}
  \item $r$ = coeficiente de correlación de Pearson
  \item $s_Y$ = desviación estándar de $Y$
  \item $s_X$ = desviación estándar de $X$
\end{itemize}

\vspace{0.3cm}
\textbf{Interpretación:}
\begin{itemize}
  \item Si $r = 0$ $\Rightarrow$ $\hat{\beta}_1 = 0$ (no hay relación lineal)
  \item Si $s_X = s_Y$ y $r = 1$ $\Rightarrow$ $\hat{\beta}_1 = 1$ (relación 1 a 1)
  \item El signo de $\hat{\beta}_1$ es el mismo que el de $r$
\end{itemize}

\vspace{0.2cm}
\begin{exampleblock}{Nota}
Aunque $\hat{\beta}_1$ depende de $r$, la regresión aporta más información: cuantifica el cambio esperado en $Y$ por unidad de cambio en $X$.
\end{exampleblock}
\end{frame}

\begin{frame}[fragile]{Aplicación en R: \texttt{lm()}}
\begin{lstlisting}
# Ajustar modelo de regresion lineal simple
modelo1 <- lm(MOD_INGLES_PUNT ~ INGLES_SABER11_PUNT,
              data = icfes_ni)

# Ver coeficientes estimados
coef(modelo1)

#        (Intercept) INGLES_SABER11_PUNT
#              45.23                 1.82

# Interpretar
# beta_0 = 45.23: puntaje esperado en Saber Pro si Saber 11 = 0
# beta_1 = 1.82: por cada punto adicional en Saber 11,
#                esperamos 1.82 puntos mas en Saber Pro
\end{lstlisting}

\vspace{0.2cm}
\begin{alertblock}{Cuidado con la extrapolación}
$\hat{\beta}_0 = 45.23$ representa el intercepto teórico, pero en la práctica nadie obtiene 0 en Saber 11. No interprete el intercepto fuera del rango de los datos.
\end{alertblock}
\end{frame}

\begin{frame}[fragile]{Valores ajustados y residuos}
\begin{lstlisting}
# Valores ajustados (predicciones)
y_hat <- fitted(modelo1)
head(y_hat)

# 153.2 158.4 162.1 147.8 ...

# Residuos
residuos <- residuals(modelo1)
head(residuos)

# 4.8 -2.4 1.9 -3.8 ...

# Verificar: la suma de residuos es (casi) 0
sum(residuos)

# [1] -1.42e-12  (numericamente 0)

# Verificar: la correlacion entre X y residuos es (casi) 0
cor(icfes_ni$INGLES_SABER11_PUNT, residuos)

# [1] -3.5e-17  (numericamente 0)
\end{lstlisting}

\textbf{Propiedades OLS:} $\sum e_i = 0$ y $\sum X_i e_i = 0$ (ortogonalidad).
\end{frame}

\begin{frame}[fragile]{Visualización de residuos}
\begin{lstlisting}
# Grafico de residuos vs valores ajustados
plot(fitted(modelo1), residuals(modelo1),
     xlab = "Valores ajustados",
     ylab = "Residuos",
     main = "Residuos vs Valores Ajustados",
     pch = 20, col = "steelblue")
abline(h = 0, col = "red", lwd = 2, lty = 2)

# Idealmente, los residuos deberian estar dispersos
# aleatoriamente alrededor de 0, sin patron visible
\end{lstlisting}

\vspace{0.2cm}
\begin{block}{Diagnóstico}
\begin{itemize}
  \item Patrón en forma de embudo $\Rightarrow$ heterocedasticidad
  \item Patrón curvo $\Rightarrow$ no linealidad
  \item Dispersión aleatoria $\Rightarrow$ supuestos plausibles
\end{itemize}
\end{block}
\end{frame}

% ============================================================================
\section{Supuestos de Gauss-Markov}
% ============================================================================

\begin{frame}{Los supuestos clásicos de regresión lineal}
Para que OLS sea el \textbf{Mejor Estimador Lineal Insesgado} (BLUE: Best Linear Unbiased Estimator), se requieren los siguientes supuestos:

\vspace{0.2cm}
\begin{enumerate}
  \item \textbf{Linealidad:} $E(Y|X) = \beta_0 + \beta_1 X$\\
    La relación entre $X$ e $Y$ es lineal en los parámetros.

  \item \textbf{Independencia:} Las observaciones son independientes entre sí.

  \item \textbf{Homocedasticidad:} $\text{Var}(\varepsilon_i | X_i) = \sigma^2$ (constante)\\
    La varianza del error es constante para todos los valores de $X$.

  \item \textbf{Normalidad:} $\varepsilon_i \sim N(0, \sigma^2)$\\
    Los errores se distribuyen normalmente (necesario para inferencia).

  \item \textbf{Exogeneidad:} $E(\varepsilon_i | X_i) = 0$\\
    El error no está correlacionado con $X$ (no hay sesgo de endogeneidad).
\end{enumerate}
\end{frame}

\begin{frame}{Consecuencias de violar los supuestos}
\begin{alertblock}{Violación de linealidad}
$\Rightarrow$ Estimadores sesgados. Solución: transformar variables (log, cuadrática, etc.) o usar modelos no lineales.
\end{alertblock}

\begin{alertblock}{Violación de independencia}
$\Rightarrow$ Errores estándar incorrectos (muy optimistas). Solución: errores estándar robustos a cluster.
\end{alertblock}

\begin{alertblock}{Violación de homocedasticidad}
$\Rightarrow$ Estimadores no eficientes, errores estándar sesgados. Solución: errores estándar robustos (HC) o WLS.
\end{alertblock}

\begin{alertblock}{Violación de normalidad}
$\Rightarrow$ Inferencia (pruebas t, IC) no válida en muestras pequeñas. Solución: bootstrap, transformaciones, o confiar en TCL si $n$ es grande.
\end{alertblock}
\end{frame}

\begin{frame}[fragile]{Verificación de supuestos: gráficos diagnósticos}
\begin{lstlisting}
# R proporciona 4 graficos diagnosticos automaticos
par(mfrow = c(2, 2))  # 2x2 panel
plot(modelo1)
par(mfrow = c(1, 1))  # restaurar

# 1. Residuals vs Fitted: verifica linealidad y homocedasticidad
# 2. Q-Q Plot: verifica normalidad de residuos
# 3. Scale-Location: verifica homocedasticidad
# 4. Residuals vs Leverage: identifica observaciones influyentes
\end{lstlisting}

\vspace{0.2cm}
\textbf{Lectura:}
\begin{itemize}
  \item Residuals vs Fitted: línea roja horizontal y dispersión uniforme $\Rightarrow$ OK
  \item Q-Q Plot: puntos sobre la diagonal $\Rightarrow$ normalidad
  \item Scale-Location: línea horizontal $\Rightarrow$ homocedasticidad
  \item Residuals vs Leverage: puntos fuera de las bandas de Cook $\Rightarrow$ outliers influyentes
\end{itemize}
\end{frame}

% ============================================================================
\section{Coeficiente de determinación $R^2$}
% ============================================================================

\begin{frame}{Descomposición de la variabilidad: SST = SSR + SSE}
La variabilidad total de $Y$ se descompone en dos partes:

\begin{block}{Sumas de cuadrados}
\begin{itemize}
  \item \textbf{SST} (Total Sum of Squares): $\sum (Y_i - \bar{Y})^2$\\
    Variabilidad total de $Y$ respecto a su media.

  \item \textbf{SSR} (Regression Sum of Squares): $\sum (\hat{Y}_i - \bar{Y})^2$\\
    Variabilidad \textbf{explicada} por el modelo.

  \item \textbf{SSE} (Error Sum of Squares): $\sum (Y_i - \hat{Y}_i)^2 = \sum e_i^2$\\
    Variabilidad \textbf{no explicada} (residual).
\end{itemize}
\end{block}

\vspace{0.2cm}
\textbf{Relación fundamental:}
\[
\text{SST} = \text{SSR} + \text{SSE}
\]

Dividiendo por SST:
\[
1 = \frac{\text{SSR}}{\text{SST}} + \frac{\text{SSE}}{\text{SST}}
\]
\end{frame}

\begin{frame}{Coeficiente de determinación $R^2$}
\begin{block}{Definición}
\[
R^2 = \frac{\text{SSR}}{\text{SST}} = 1 - \frac{\text{SSE}}{\text{SST}}
\]
\end{block}

\vspace{0.2cm}
\textbf{Interpretación:}
\begin{itemize}
  \item $R^2$ = proporción de la variabilidad de $Y$ \textbf{explicada} por $X$
  \item Rango: $R^2 \in [0, 1]$
  \item $R^2 = 0$ $\Rightarrow$ $X$ no explica nada de $Y$ (SSR = 0)
  \item $R^2 = 1$ $\Rightarrow$ ajuste perfecto (SSE = 0, todos los puntos sobre la línea)
\end{itemize}

\vspace{0.3cm}
\begin{exampleblock}{Ejemplo}
Si $R^2 = 0.65$, decimos que el 65\% de la variabilidad del puntaje de inglés en Saber Pro se explica por el puntaje en Saber 11. El 35\% restante se debe a otros factores no incluidos en el modelo.
\end{exampleblock}
\end{frame}

\begin{frame}{Relación entre $R^2$ y correlación}
En regresión lineal simple (un solo predictor), se cumple:
\[
R^2 = r^2
\]
donde $r$ es el coeficiente de correlación de Pearson entre $X$ e $Y$.

\vspace{0.3cm}
\textbf{Ejemplo:}
\begin{itemize}
  \item Si $r = 0.8$ $\Rightarrow$ $R^2 = 0.64$ (64\% de variabilidad explicada)
  \item Si $r = -0.5$ $\Rightarrow$ $R^2 = 0.25$ (25\% de variabilidad explicada)
\end{itemize}

\vspace{0.3cm}
\begin{alertblock}{Nota importante}
$R^2$ es siempre positivo (es un cuadrado), pero $r$ puede ser negativo. El signo de la relación se ve en $\hat{\beta}_1$ (o en $r$), no en $R^2$.
\end{alertblock}
\end{frame}

\begin{frame}[fragile]{Calculando $R^2$ en R}
\begin{lstlisting}
# Obtener R^2 del modelo
summary(modelo1)$r.squared

# [1] 0.6524

# Verificar con la correlacion
r <- cor(icfes_ni$MOD_INGLES_PUNT,
         icfes_ni$INGLES_SABER11_PUNT,
         use = "complete.obs")
r^2

# [1] 0.6524

# Calcular manualmente SST, SSR, SSE
SST <- sum((icfes_ni$MOD_INGLES_PUNT - mean(icfes_ni$MOD_INGLES_PUNT))^2)
SSE <- sum(residuals(modelo1)^2)
SSR <- SST - SSE

R2_manual <- SSR / SST
R2_manual

# [1] 0.6524
\end{lstlisting}
\end{frame}

\begin{frame}{Precauciones al interpretar $R^2$}
\begin{alertblock}{$R^2$ alto $\not\Rightarrow$ causalidad}
Una correlación alta no implica que $X$ \textbf{cause} $Y$. Puede haber:
\begin{itemize}
  \item Causalidad inversa ($Y$ causa $X$)
  \item Variables omitidas (confounders)
  \item Correlación espuria
\end{itemize}
\end{alertblock}

\begin{alertblock}{$R^2$ alto $\not\Rightarrow$ buen modelo de predicción}
Un modelo puede tener $R^2$ alto en la muestra pero predecir mal fuera de la muestra (overfitting).
\end{alertblock}

\begin{alertblock}{$R^2$ bajo $\not\Rightarrow$ modelo inútil}
Incluso con $R^2$ bajo, el modelo puede revelar una relación importante y estadísticamente significativa. En ciencias sociales, $R^2 = 0.20$ puede ser razonable.
\end{alertblock}
\end{frame}

% ============================================================================
\section{Pruebas de significancia}
% ============================================================================

\begin{frame}{Varianza del error y MSE}
Para hacer inferencia sobre los parámetros, necesitamos estimar $\sigma^2 = \text{Var}(\varepsilon_i)$.

\begin{block}{Estimador de $\sigma^2$}
\[
s^2 = \text{MSE} = \frac{\text{SSE}}{n - 2}
\]
donde $n - 2$ son los grados de libertad ($n$ observaciones menos 2 parámetros estimados).
\end{block}

\vspace{0.2cm}
\textbf{MSE} = Mean Squared Error (cuadrado medio del error residual).

\vspace{0.2cm}
La raíz cuadrada de MSE se llama \textbf{error estándar de la regresión}:
\[
s = \sqrt{\text{MSE}} = \text{residual standard error}
\]

\vspace{0.2cm}
\begin{exampleblock}{Interpretación}
$s$ representa la desviación estándar típica de los residuos. Si $s = 12$, en promedio los valores observados se desvían $\pm 12$ puntos de la línea de regresión.
\end{exampleblock}
\end{frame}

\begin{frame}{Error estándar de $\hat{\beta}_1$}
El estimador $\hat{\beta}_1$ tiene una distribución muestral con:
\[
\text{SE}(\hat{\beta}_1) = \frac{s}{\sqrt{\sum(X_i - \bar{X})^2}} = \frac{s}{\sqrt{S_{XX}}}
\]

\vspace{0.2cm}
\textbf{Interpretación:}
\begin{itemize}
  \item Mayor variabilidad en $X$ (mayor $S_{XX}$) $\Rightarrow$ menor error estándar (estimación más precisa)
  \item Mayor variabilidad residual ($s$ alto) $\Rightarrow$ mayor error estándar (estimación menos precisa)
\end{itemize}

\vspace{0.3cm}
\begin{block}{Distribución bajo los supuestos}
Si los errores son normales:
\[
\frac{\hat{\beta}_1 - \beta_1}{\text{SE}(\hat{\beta}_1)} \sim t_{n-2}
\]
\end{block}
\end{frame}

\begin{frame}{Prueba t para la pendiente}
\begin{block}{Hipótesis más común}
\begin{itemize}
  \item $H_0: \beta_1 = 0$ (no hay relación lineal entre $X$ e $Y$)
  \item $H_1: \beta_1 \neq 0$ (existe relación lineal)
\end{itemize}
\end{block}

\vspace{0.2cm}
\textbf{Estadístico de prueba:}
\[
t = \frac{\hat{\beta}_1}{\text{SE}(\hat{\beta}_1)} \sim t_{n-2} \quad \text{bajo } H_0
\]

\vspace{0.2cm}
\textbf{Regla de decisión:}
\begin{itemize}
  \item Rechazamos $H_0$ si $|t| > t_{\alpha/2, n-2}$
  \item O equivalentemente, si p-valor $< \alpha$
\end{itemize}

\vspace{0.2cm}
\begin{exampleblock}{Interpretación}
Si rechazamos $H_0$, concluimos que la pendiente es \textbf{estadísticamente diferente de cero}, es decir, existe una relación lineal significativa entre $X$ e $Y$.
\end{exampleblock}
\end{frame}

\begin{frame}{Intervalo de confianza para $\beta_1$}
Un IC del $(1-\alpha) \times 100\%$ para $\beta_1$ es:
\[
\hat{\beta}_1 \pm t_{\alpha/2, n-2} \times \text{SE}(\hat{\beta}_1)
\]

\vspace{0.2cm}
\textbf{Interpretación:}
\begin{itemize}
  \item Con 95\% de confianza, el verdadero valor de $\beta_1$ está en este intervalo.
  \item Si el intervalo \textbf{no contiene 0}, la pendiente es significativa al nivel $\alpha$.
\end{itemize}

\vspace{0.3cm}
\begin{exampleblock}{Ejemplo}
IC 95\% para $\beta_1$: $[1.65, 1.99]$
\begin{itemize}
  \item No contiene 0 $\Rightarrow$ la relación es significativa
  \item Interpretación: con 95\% de confianza, por cada punto adicional en Saber 11, el puntaje en Saber Pro aumenta entre 1.65 y 1.99 puntos
\end{itemize}
\end{exampleblock}
\end{frame}

\begin{frame}[fragile]{Aplicación en R: \texttt{summary(lm())}}
\begin{lstlisting}
# Resumen completo del modelo
summary(modelo1)

# Call:
# lm(formula = MOD_INGLES_PUNT ~ INGLES_SABER11_PUNT, data = icfes_ni)
#
# Residuals:
#     Min      1Q  Median      3Q     Max
# -45.231  -8.124   0.342   8.567  39.128
#
# Coefficients:
#                      Estimate Std. Error t value Pr(>|t|)
# (Intercept)           45.2314     1.2345   36.64   <2e-16 ***
# INGLES_SABER11_PUNT    1.8234     0.0198   92.09   <2e-16 ***
# ---
# Signif. codes:  0 '***' 0.001 '**' 0.01 '*' 0.05 '.' 0.1 ' ' 1
#
# Residual standard error: 12.35 on 5083 degrees of freedom
# Multiple R-squared:  0.6524, Adjusted R-squared:  0.6523
# F-statistic: 8481 on 1 and 5083 DF,  p-value: < 2.2e-16
\end{lstlisting}
\end{frame}

\begin{frame}{Interpretando el output de \texttt{summary()}}
\textbf{Coefficients:}
\begin{itemize}
  \item \texttt{Estimate}: $\hat{\beta}_0 = 45.23$, $\hat{\beta}_1 = 1.82$
  \item \texttt{Std. Error}: SE($\hat{\beta}_0$) = 1.23, SE($\hat{\beta}_1$) = 0.020
  \item \texttt{t value}: $t = \hat{\beta}/\text{SE}(\hat{\beta})$ = 92.09 para $\beta_1$
  \item \texttt{Pr(>|t|)}: p-valor $< 2 \times 10^{-16}$ (altamente significativo)
\end{itemize}

\vspace{0.2cm}
\textbf{Residual standard error:} $s = 12.35$ (desviación típica de residuos)

\vspace{0.2cm}
\textbf{Multiple R-squared:} $R^2 = 0.6524$ (65.24\% de variabilidad explicada)

\vspace{0.2cm}
\textbf{F-statistic:} $F = 8481$ con p-valor $< 2.2 \times 10^{-16}$ (prueba global del modelo)

\vspace{0.2cm}
\begin{alertblock}{Conclusión}
Existe una relación lineal \textbf{muy significativa} entre el inglés de Saber 11 y Saber Pro. Por cada punto adicional en Saber 11, esperamos 1.82 puntos más en Saber Pro.
\end{alertblock}
\end{frame}

\begin{frame}[fragile]{Intervalos de confianza para los coeficientes}
\begin{lstlisting}
# IC 95% para todos los coeficientes
confint(modelo1, level = 0.95)

#                          2.5 %    97.5 %
# (Intercept)           42.81136  47.65144
# INGLES_SABER11_PUNT    1.78455   1.86225

# IC 99% (mas amplio)
confint(modelo1, level = 0.99)

#                          0.5 %    99.5 %
# (Intercept)           42.05312  48.40968
# INGLES_SABER11_PUNT    1.77345   1.87335
\end{lstlisting}

\vspace{0.2cm}
\textbf{Interpretación:}
\begin{itemize}
  \item Con 95\% de confianza, $\beta_1 \in [1.78, 1.86]$
  \item El intervalo no contiene 0 $\Rightarrow$ la relación es significativa
\end{itemize}
\end{frame}

\begin{frame}{Prueba F global}
En regresión simple, la \textbf{prueba F global} es equivalente a la prueba t para $\beta_1$:

\begin{block}{Hipótesis}
\begin{itemize}
  \item $H_0$: $\beta_1 = 0$ (el modelo no es útil)
  \item $H_1$: $\beta_1 \neq 0$ (el modelo es útil)
\end{itemize}
\end{block}

\vspace{0.2cm}
\textbf{Estadístico:}
\[
F = \frac{\text{MSR}}{\text{MSE}} = \frac{\text{SSR}/1}{\text{SSE}/(n-2)} \sim F_{1, n-2} \quad \text{bajo } H_0
\]

\vspace{0.2cm}
\textbf{Relación con prueba t:}
\[
F = t^2
\]

En regresión simple, $F$ y $t$ dan la misma conclusión. En regresión múltiple, $F$ prueba si \textbf{todos} los coeficientes son cero simultáneamente (veremos en S10).
\end{frame}

% ============================================================================
\section{Modelo con variable dummy}
% ============================================================================

\begin{frame}{Variables dummy (indicadoras)}
Una \textbf{variable dummy} es una variable binaria que toma valores 0 o 1 para representar categorías.

\vspace{0.2cm}
\begin{exampleblock}{Ejemplo}
Queremos comparar el puntaje de inglés entre IES privadas y públicas:
\[
\text{PRIVADA} = \begin{cases}
1 & \text{si IES privada} \\
0 & \text{si IES pública}
\end{cases}
\]
\end{exampleblock}

\vspace{0.2cm}
\textbf{Modelo:}
\[
\text{MOD\_INGLES\_PUNT}_i = \beta_0 + \beta_1 \times \text{PRIVADA}_i + \varepsilon_i
\]

\vspace{0.2cm}
\textbf{Interpretación:}
\begin{itemize}
  \item $\beta_0$ = puntaje promedio en IES públicas (cuando PRIVADA = 0)
  \item $\beta_0 + \beta_1$ = puntaje promedio en IES privadas (cuando PRIVADA = 1)
  \item $\beta_1$ = diferencia promedio entre IES privadas y públicas
\end{itemize}
\end{frame}

\begin{frame}[fragile]{Regresión con dummy en R}
\begin{lstlisting}
# Crear variable dummy
icfes_ni$PRIVADA <- ifelse(icfes_ni$TIPO_IES == "privada", 1, 0)

# Ajustar modelo
modelo_dummy <- lm(MOD_INGLES_PUNT ~ PRIVADA, data = icfes_ni)

summary(modelo_dummy)

# Coefficients:
#             Estimate Std. Error t value Pr(>|t|)
# (Intercept)  154.82       0.42  368.62   <2e-16 ***
# PRIVADA        7.35       0.58   12.67   <2e-16 ***
#
# Residual standard error: 20.23 on 5083 DF
# Multiple R-squared:  0.0306, Adjusted R-squared:  0.0304
\end{lstlisting}

\vspace{0.2cm}
\textbf{Interpretación:}
\begin{itemize}
  \item Puntaje promedio en IES públicas: $\hat{\beta}_0 = 154.82$
  \item Puntaje promedio en IES privadas: $154.82 + 7.35 = 162.17$
  \item Diferencia: $\hat{\beta}_1 = 7.35$ puntos (significativa, p $< 0.001$)
\end{itemize}
\end{frame}

\begin{frame}[fragile]{Equivalencia con prueba t de dos muestras}
La regresión con una dummy es equivalente a una prueba t de dos muestras independientes:

\begin{lstlisting}
# Prueba t tradicional
t.test(MOD_INGLES_PUNT ~ TIPO_IES, data = icfes_ni,
       var.equal = TRUE)

# Two Sample t-test
#
# data:  MOD_INGLES_PUNT by TIPO_IES
# t = -12.67, df = 5083, p-value < 2.2e-16
# 95 percent confidence interval:
#  -8.486 -6.214
# sample estimates:
# mean in group privada mean in group publica
#              162.17              154.82
\end{lstlisting}

\vspace{0.2cm}
\begin{block}{Observación}
El estadístico $t$ de la regresión ($t = 12.67$) es idéntico (en valor absoluto) al de la prueba t de dos muestras. La regresión es una generalización poderosa.
\end{block}
\end{frame}

\begin{frame}[fragile]{Visualización: boxplot por grupo}
\begin{lstlisting}
# Comparar graficamente
library(ggplot2)

ggplot(icfes_ni, aes(x = TIPO_IES, y = MOD_INGLES_PUNT,
                     fill = TIPO_IES)) +
  geom_boxplot(alpha = 0.7) +
  geom_hline(yintercept = coef(modelo_dummy)[1],
             linetype = "dashed", color = "blue", size = 1) +
  geom_hline(yintercept = coef(modelo_dummy)[1] + coef(modelo_dummy)[2],
             linetype = "dashed", color = "red", size = 1) +
  labs(title = "Puntaje de ingles por tipo de IES",
       x = "Tipo de IES",
       y = "MOD_INGLES_PUNT") +
  scale_fill_manual(values = c("steelblue", "coral")) +
  theme_minimal()
\end{lstlisting}

Las líneas punteadas muestran los valores predichos por el modelo ($\hat{\beta}_0$ y $\hat{\beta}_0 + \hat{\beta}_1$).
\end{frame}

% ============================================================================
\section{Ejemplo completo en R}
% ============================================================================

\begin{frame}[fragile]{Ejemplo integrador: paso a paso}
\textbf{Pregunta:} ¿Cómo se relaciona el puntaje cuantitativo en Saber Pro (MOD\_RAZONA\_CUANT\_PUNT) con el puntaje de inglés en Saber 11?

\begin{lstlisting}
# 1. Exploracion inicial
plot(icfes_ni$INGLES_SABER11_PUNT,
     icfes_ni$MOD_RAZONA_CUANT_PUNT,
     xlab = "Ingles Saber 11",
     ylab = "Razonamiento Cuantitativo Saber Pro",
     main = "Relacion entre ingles y razonamiento cuantitativo",
     pch = 20, col = rgb(0, 0, 1, 0.3))

# 2. Calcular correlacion
cor(icfes_ni$INGLES_SABER11_PUNT,
    icfes_ni$MOD_RAZONA_CUANT_PUNT,
    use = "complete.obs")

# [1] 0.58  (correlacion moderada-alta)
\end{lstlisting}
\end{frame}

\begin{frame}[fragile]{Ejemplo integrador: ajuste del modelo}
\begin{lstlisting}
# 3. Ajustar modelo de regresion
modelo2 <- lm(MOD_RAZONA_CUANT_PUNT ~ INGLES_SABER11_PUNT,
              data = icfes_ni)

# 4. Ver resumen
summary(modelo2)

# Coefficients:
#                      Estimate Std. Error t value Pr(>|t|)
# (Intercept)           98.456      1.834   53.68   <2e-16 ***
# INGLES_SABER11_PUNT    0.987      0.029   34.03   <2e-16 ***
#
# Residual standard error: 18.36 on 5083 DF
# Multiple R-squared:  0.3364, Adjusted R-squared:  0.3363

# 5. IC 95% para los coeficientes
confint(modelo2)

#                          2.5 %   97.5 %
# (Intercept)           94.8604  102.051
# INGLES_SABER11_PUNT    0.9296    1.044
\end{lstlisting}
\end{frame}

\begin{frame}[fragile]{Ejemplo integrador: diagnóstico}
\begin{lstlisting}
# 6. Graficos diagnosticos
par(mfrow = c(2, 2))
plot(modelo2)
par(mfrow = c(1, 1))

# 7. Verificar supuestos
# Normalidad de residuos
shapiro.test(sample(residuals(modelo2), 5000))

# W = 0.9967, p-value = 0.0002
# Leve desviacion de normalidad, pero con n grande
# el TCL asegura validez asintotica

# Homocedasticidad (test de Breusch-Pagan)
library(lmtest)
bptest(modelo2)

# BP = 23.45, p-value = 1.3e-06
# Hay evidencia de heterocedasticidad
# (varianza no constante)
\end{lstlisting}
\end{frame}

\begin{frame}[fragile]{Ejemplo integrador: predicción}
\begin{lstlisting}
# 8. Hacer predicciones
# Prediccion puntual para Saber 11 = 70
nuevo_dato <- data.frame(INGLES_SABER11_PUNT = 70)
predict(modelo2, newdata = nuevo_dato)

# [1] 167.55

# Prediccion con IC 95%
predict(modelo2, newdata = nuevo_dato, interval = "confidence")

#      fit    lwr    upr
# 1 167.55 166.23 168.87

# Prediccion con IP 95% (prediction interval)
predict(modelo2, newdata = nuevo_dato, interval = "prediction")

#      fit     lwr     upr
# 1 167.55 131.42 203.68
\end{lstlisting}

\textbf{Diferencia:} IC para la \textbf{media} vs IP para una \textbf{observación individual} (IP siempre más amplio).
\end{frame}

\begin{frame}[fragile]{Ejemplo integrador: visualización final}
\begin{lstlisting}
# 9. Grafico con linea de regresion e intervalos
library(ggplot2)

ggplot(icfes_ni, aes(x = INGLES_SABER11_PUNT,
                     y = MOD_RAZONA_CUANT_PUNT)) +
  geom_point(alpha = 0.2, color = "steelblue") +
  geom_smooth(method = "lm", se = TRUE, color = "red",
              fill = "pink", alpha = 0.3) +
  labs(title = "Ingles (Saber 11) vs Razonamiento Cuantitativo (Saber Pro)",
       subtitle = paste("R^2 = 0.336, p < 0.001"),
       x = "Puntaje Ingles Saber 11",
       y = "Puntaje Razonamiento Cuantitativo Saber Pro") +
  theme_minimal() +
  theme(plot.title = element_text(hjust = 0.5),
        plot.subtitle = element_text(hjust = 0.5))
\end{lstlisting}
\end{frame}

\begin{frame}{Resumen de la sesión}
\begin{block}{Conceptos clave}
\begin{itemize}
  \item Modelo de regresión lineal simple: $Y = \beta_0 + \beta_1 X + \varepsilon$
  \item OLS minimiza $\sum e_i^2$ para estimar $\hat{\beta}_0$ y $\hat{\beta}_1$
  \item La línea OLS pasa por $(\bar{X}, \bar{Y})$
  \item $R^2$ mide la proporción de variabilidad explicada
  \item Supuestos de Gauss-Markov: linealidad, independencia, homocedasticidad, normalidad, exogeneidad
\end{itemize}
\end{block}

\begin{block}{Inferencia}
\begin{itemize}
  \item Prueba t para $H_0: \beta_1 = 0$
  \item IC para $\beta_1$: $\hat{\beta}_1 \pm t_{\alpha/2, n-2} \times \text{SE}(\hat{\beta}_1)$
  \item Regresión con dummy $\equiv$ t-test de dos muestras
\end{itemize}
\end{block}

\begin{exampleblock}{Próxima sesión}
Regresión lineal múltiple: incorporar varios predictores simultáneamente.
\end{exampleblock}
\end{frame}

\end{document}
